% Approximately 500 words or less; try not to use formulas or special characters
% If you don't want an initial indentation, do \noindent at the start of the abstract

Despite advances in mobility technology, cities still struggle to design bus systems that are efficient, equitable, and adaptable to evolving urban needs. Computational bus service design aims to identify bus routes which match historic ridership demand while complying with requirements from the built environment, expected reliability, equity, efficiency, and environmental impact, but very little work has been done on integrating equity analyses into these computational network optimizations.  In this thesis, I explore how to conduct bus network redesigns in a method that balances service quality and cost with equity and accessibility factors. This is done by creating a blueprint using an applied category theory framework that can take into account transit network data and capture the complex interplay of various stakeholders. This formalism will assist city planners, transit operators, and passengers, and is applied as a case study to a real-world transit agency in Framingham, MA to assess feasibility and accuracy, as well as to identify areas for further improvement and research. 